\section{System Design}
\label{sec:system}

\textsc{Pichay} is a transparent HTTP proxy that interposes between
an agentic AI client and the inference API. On each request, it
receives the full message array assembled by the client, applies
context management policies, and forwards the modified request to
the inference provider. The client is unaware of the interposition.

\subsection{Architecture}

The proxy operates at the Messages API level. It receives JSON
requests containing the system prompt, tool definitions, and message
history. It returns the inference provider's response unmodified.
The interposition point is the message array: the proxy can inspect,
measure, modify, or replace any message before forwarding.

This design has three properties: (1)~it requires no changes to the
client, the model, or the inference API; (2)~the client retains the
full, unmodified conversation history, providing a backing store
from which evicted content can be faulted back in; and (3)~the proxy
sees the complete request on every turn, enabling stateful
management without persistent connections.

\subsection{Garbage Collection vs.\ Paging}

Not all tool results are equal. We distinguish two categories:

\textbf{Garbage collection} applies to ephemeral tool outputs---Bash
command results, search outputs, directory listings---that have no
stable identity. Once consumed, they cannot be meaningfully
re-requested. Removing them is reclamation of dead content.

\textbf{Paging} applies to addressable content---file reads, plan
documents, specifications---that have stable identity (a file path).
Removing these creates a fault risk: the model may need the content
again and must re-read the file.

The distinction matters for fault rate calculation. Only paging
(Read evictions) can produce faults. Garbage collection cannot.
Conflating them inflates the eviction denominator and deflates the
apparent fault rate.

\subsection{Eviction Policy}

The current policy is FIFO by user-turn age: tool results older than
$\tau$~user turns and larger than $s_{\min}$~bytes are candidates
for eviction ($\tau = 4$, $s_{\min} = 500$ in all experiments). This
is deliberately simple. The research question is how far a minimal
policy extends before requiring sophistication.

Evicted content is replaced with a retrieval handle:

\begin{quote}
\small\texttt{[Paged out: Read /path/to/file.py (8,192 bytes,
187 lines). Re-read the file if you need its content.]}
\end{quote}

The handle serves three functions: (1)~it identifies what was
removed; (2)~it states how to recover the content; and (3)~it
occupies substantially less context than the original.

\subsection{Page Fault Detection}

When the model issues a tool call that matches an evicted entry (same
tool name and arguments), the proxy detects a \emph{page fault}: the
model is requesting content it previously had but lost to eviction.
The evicted entry is indexed by a key derived from the tool name and
input arguments (e.g., file path for Read operations).

Page faults are observable events. They measure the cost of eviction:
each fault is a wasted tool call round-trip that would not have
occurred if the content had remained resident. At inference-time
pricing, faults have direct monetary cost.

\subsection{Fault-Driven Pinning}
\label{sec:pinning}

The simplest upgrade to FIFO: if evicting a page caused a fault,
don't evict it again. One fault per file, permanently pinned for the
session. The algorithm:

\begin{enumerate}[nosep]
  \item On eviction of addressable content: record the file path and
    content hash.
  \item On page fault: record the evicted content's hash in a fault
    history table.
  \item On next eviction attempt for the same path: if the current
    content hash matches the fault history entry, pin the page---skip
    eviction permanently.
  \item On new read of a pinned path with different content (file was
    edited): unpin. The old version is stale; the new version starts
    a fresh fault cycle.
\end{enumerate}

The content hash comparison prevents false pins: if a file changed
between eviction and re-read, the eviction was correct (stale data
was removed). Pinning only applies when the model demonstrably
needed exactly what was taken away. The unpin-on-edit rule handles
the edit/review cycle common in coding sessions.

\subsection{Retrieval Handles as Anchors}

The eviction summary was designed as a space-saving marker. In
practice, it functions as a late-binding retrieval handle---an anchor
that stores minimal metadata and resolves to full content on demand.

The handle resolves to \emph{current} content, not the original
evicted content. A file edited since eviction materializes at its
new state when faulted in. This temporal property is a feature: the
model always gets the latest version, and stale cached content is
impossible.

Behavioral evidence: when a fresh model instance resumed a session
containing paged-out content, it unprompted stated: ``Let me re-read
the files I need since some were paged out.'' The model recognized
the handles, understood what was missing, and chose to fault content
in before acting. The handle's format carries its own semantics---no
instruction was needed.

\subsection{Cooperative Memory Management}
\label{sec:cooperative}

In OS virtual memory, the application is non-cooperative. It never
voluntarily releases pages. The OS must infer the working set from
access patterns because the CPU will not say what it needs next.
Decades of replacement algorithms---FIFO, LRU, Clock, Working
Set---exist because of this fundamental asymmetry.

LLMs break this assumption. The model has \emph{incentive} to
cooperate with the memory manager: cleaner context means better
attention, better output quality, and longer session life. The
processor wants to help manage its own cache. This changes the
design space.

\textsc{Pichay} exploits this through \emph{phantom tools}---tool
definitions injected by the proxy that the framework never sees.
The model can call them; the proxy intercepts the calls from the
streaming response before the framework receives them, handles
them internally, and injects coherent tool results on the next
turn. The framework is unaware of the side channel.

Two phantom tools are defined:

\begin{description}[nosep]
  \item[\texttt{memory\_release(paths)}] The model signals it no
    longer needs specific files. The proxy marks them for immediate
    eviction, bypassing the age threshold. This is the reference
    bit---provided voluntarily by the processor.
  \item[\texttt{memory\_fault(paths)}] The model requests evicted
    content restored from the proxy's backing store. The proxy
    resolves the fault from its eviction cache without a file system
    round trip---faster and cheaper than a real Read tool call.
\end{description}

The second cooperative channel is \emph{cleanup tags}---structured
directives embedded in the model's output text, parsed by the proxy
before forwarding to the framework. Where phantom tools are
proxy-to-model (the proxy offers capabilities), cleanup tags are
model-to-proxy (the model initiates management). Four operations are
defined:

\begin{description}[nosep]
  \item[\texttt{drop: block:ID}] Immediately evict a specific block.
  \item[\texttt{summarize: block:ID "text"}] Replace a block with a
    compact summary---lossy compression authored by the model, not
    the system.
  \item[\texttt{anchor: block:ID}] Pin a block against eviction.
  \item[\texttt{collapse: turns N-M "text"}] Compress a range of
    conversation turns into a single synthetic summary block,
    removing all intermediate content.
\end{description}

The collapse operation is particularly significant: it enables L3
compaction (Section~\ref{sec:l3}) initiated by the model itself,
based on its understanding of which dialogue was scaffolding and
which produced durable outcomes.

The key insight: in hardware memory hierarchies, the application is
adversarial or indifferent. The entire replacement algorithm
literature assumes non-cooperation. Cooperative demand paging---where
the processor voluntarily releases cold pages and explicitly
requests faults---is a new point in the design space, enabled by the
fact that LLMs experience degraded output quality under context
pressure and can learn that cooperation improves their own
performance. The two channels (phantom tools and cleanup tags)
provide both directions of the cooperative protocol: the proxy
informs the model of memory state, and the model directs the proxy
to act on its cognitive assessment of what is cold.

\subsection{Graduated Pressure Zones}
\label{sec:pressure}

The eviction policy must respond to context pressure, but a binary
threshold (evict/don't-evict) is too coarse. \textsc{Pichay} defines
four pressure zones based on token consumption:

\begin{description}[nosep]
  \item[Normal] ($< 60\text{K}$ tokens): No intervention. The proxy
    observes and logs.
  \item[Advisory] ($60\text{K}$--$100\text{K}$): The proxy injects
    memory pressure information into the model's context---current
    fill percentage, the five largest resident blocks, and available
    cleanup operations. The model can act cooperatively or ignore the
    advisory. This is the graduated equivalent of the ``low memory''
    notification in desktop operating systems.
  \item[Involuntary] ($100\text{K}$--$120\text{K}$): The proxy begins
    automatic eviction using the FIFO age policy. The model is
    informed but not consulted.
  \item[Aggressive] ($\geq 120\text{K}$): Emergency eviction with
    relaxed thresholds. Context survival takes priority over working
    set preservation.
\end{description}

The advisory zone is the cooperative innovation. Rather than
evicting silently (the OS approach) or crashing at capacity (the
current agentic AI approach), the system provides the model with
enough information to make intelligent cleanup decisions. The
60K threshold was chosen to provide approximately 40K tokens of
runway before involuntary eviction---enough for the model to
complete a coherent thought and emit cleanup tags before losing
agency over the process.

\subsection{L3: Rolling Conversation Compaction}
\label{sec:l3}

The proxy manages tool results (L1/L2). Non-tool content---conversation
history, planning dialogue, orientation reads---is managed by the
collapse operation (Section~\ref{sec:cooperative}). This is L3 in the
memory hierarchy: session history compressed with declared losses into
structured summaries.

The collapse operation \texttt{collapse: turns N-M "summary"} replaces
all blocks in a contiguous turn range with a single synthetic block
containing the model-authored summary. The original content is not
stored---this is lossy compression by design. The summary captures
\emph{outcomes} (what was decided, what was built, what failed) rather
than \emph{process} (which files were read, what intermediate steps
occurred).

Block state persists across session restarts via atomic checkpointing.
On each cleanup tag processing pass, the proxy serializes block metadata
(IDs, content hashes, sizes, turns, roles, status, summaries) to a JSON
checkpoint file using atomic write (tmp file + rename). On session
creation, the checkpoint is loaded and block tracking resumes from the
persisted state. Original content is lazily repopulated as the message
array is reprocessed.

The checkpoint is metadata-only---content is not persisted (the client's
message array is the backing store). This keeps checkpoint files small
(kilobytes, not megabytes) and avoids the consistency hazard of
maintaining two copies of message content.

% ====================================================================
